%% toc: P3 33 33.1,33.2,33.3,33.4
\section{Proofs for further material from Morley's \emph{Categoricity in Power}}
\label{sec:annex-p03-w03-p03-033}

The proofs below work in Morley's normalized setting
(\cref{def:p03-morley-standing}). Two facts from
\cref{thm:p03-morley-normalization,thm:p03-morley-types} are used
throughout: every formula with parameters from a system is equivalent
modulo \(T\) to an atomic relation applied to the variable and the
parameters, so a monomorphism preserves and reflects every formula; and
every point of \(S(A)\) is realized in an extension of \(A\), all points
at once in one extension, which may be taken to be a model because every
member of \(\mathcal N(T)\) is a subsystem of a model. For monomorphisms
\(f:A\to B\) and \(g:B\to C\), substitution of parameters is functorial,
so \((g\circ f)^*=f^*\circ g^*\).

\begin{annexitem}{annex:p03-w03-p03-033-map-category}{thm:p03-w03-p03-033-map-category}{Concrete categories and type spaces}
\begin{proof}
\emph{Clause 1.} The product \((g:B'\to C)\cdot(f:A\to B)\) is defined
exactly when \(B'=B\), and then \((gf:A\to C)\) belongs to \(\mathscr C\)
by condition (ii). Let \(f:A\to B\), \(g:B\to C\) and \(h:C\to D\) be
maps of \(\mathscr C\). Both \(h\cdot(g\cdot f)\) and \((h\cdot g)\cdot f\)
have domain \(A\) and co-domain \(D\), and their functions send every
\(x\in|A|\) to \(h(g(f(x)))\); so they are the same triple. By condition
(i), \(1_A=(\mathrm{id}_{|A|}:A\to A)\) belongs to \(\mathscr C\), and for
\(f:A\to B\) the triples \(f\cdot1_A\) and \(1_B\cdot f\) have domain
\(A\), co-domain \(B\) and function \(f\), so both equal \(f\).

\emph{Clause 2.} Members of \(\mathcal N(T)\) are relational systems,
hence mathematical objects. The identity of \(|A|\) is one-one and
preserves and reflects every relation, so it is a monomorphism. If
\(f:A\to B\) and \(g:B\to C\) are monomorphisms, \(gf\) is one-one as a
composite of one-one maps, and for every relation symbol \(R\) and tuple
\(\bar x\) from \(|A|\), \(R^A(\bar x)\) holds if and only if
\(R^B(f\bar x)\) holds, if and only if \(R^C(gf\bar x)\) holds. So \(gf\)
is a monomorphism and conditions (i) and (ii) hold.

\emph{Clause 3.} Each \(S(A)\) is a topological space
(\cref{thm:p03-morley-types}), so \(\mathscr S(T)\) is a class of
mathematical objects. For a monomorphism \(f:A\to B\), \(\tilde f\)
replaces each constant \(\bar a\) by \(\overline{f(a)}\), and
\(f^*(r)=\{\psi\in F(A):\tilde f\psi\in r\}\) for \(r\in S(B)\)
(\cref{def:p03-morley-type-space}). The map \(\tilde i_{AA}\) leaves every
formula unchanged, so \(i_{AA}^*(p)=p\) for every \(p\in S(A)\). If
\(f:A\to B\) and \(g:B\to C\) are monomorphisms, every constant of
\(L(B)\) occurring in \(\tilde f\psi\) has the form \(\overline{f(a)}\),
and \(\tilde g\) replaces it by \(\overline{g(f(a))}\); hence
\(\tilde g(\tilde f\psi)=\widetilde{gf}\psi\). For \(r\in S(C)\),
\[
 (gf)^*(r)=\{\psi:\tilde g(\tilde f\psi)\in r\}
          =\{\psi:\tilde f\psi\in g^*(r)\}
          =f^*(g^*(r)).
\]
So the identity triple of \(S(A)\) is \((i_{AA}^*:S(A)\to S(A))\), which
is in \(\mathscr C(T)\) because \(i_{AA}\) is a monomorphism, and the
composite of \((g^*:S(C)\to S(B))\) and \((f^*:S(B)\to S(A))\) is
\(((gf)^*:S(C)\to S(A))\), which is in \(\mathscr C(T)\) because \(gf\) is
a monomorphism by clause~2. Every \(f^*\) is continuous and onto by
\cref{thm:p03-morley-types}.
\end{proof}
\end{annexitem}

\begin{annexitem}{annex:p03-w03-p03-033-syntactic}{thm:p03-w03-p03-033-syntactic}{The dictionary between systems and complete theories}
\begin{proof}
The language \(L\) is relational, so an open sentence of \(L(A)\) is a
Boolean combination of atomic sentences \(R(\bar a_1,\dots,\bar a_k)\)
and \(\bar a=\bar a'\).

\emph{Clause 1.} Since \(A\in\mathcal N(T)\), \(A\) is a subsystem of a
model \(M\) of \(T\). The identity map of \(|A|\) into \(|M|\) preserves
and reflects every relation and equality, so every atomic sentence of
\(L(A)\) has the same truth value in \((A,a)_{a\in|A|}\) and in
\((M,a)_{a\in|A|}\), and therefore so does every open sentence. Hence
\((M,a)_{a\in|A|}\) is a model of \(T\cup\mathscr D(A)=T(A)\), and
\(T(A)\) is consistent. For completeness, let \(\psi\) be a sentence of
\(L(A)\), let \(\bar a_1,\dots,\bar a_n\) be the distinct constants of
\(L(A)\setminus L\) occurring in it, and let \(\chi\) be the \(L\)-formula
obtained by replacing them by distinct variables \(w_1,\dots,w_n\) not
occurring in \(\psi\). If \(n=0\), \(\psi\) is a sentence of \(L\), and
\(T\), which is complete (\cref{thm:p03-morley-normalization}), proves
\(\psi\) or \(\neg\psi\). If \(n\geq1\), Lemma~1.1(e), the clause of
\cref{thm:p03-morley-normalization} on atomic equivalents, gives a
relation symbol \(R\) of degree \(n\) with
\(T\vdash\forall w_1\cdots\forall w_n(\chi\leftrightarrow
R(w_1,\dots,w_n))\), so \(T\vdash\psi\leftrightarrow
R(\bar a_1,\dots,\bar a_n)\). The open sentence
\(R(\bar a_1,\dots,\bar a_n)\) or its negation is true in
\((A,a)_{a\in|A|}\) and so belongs to \(\mathscr D(A)\). Hence
\(T(A)\vdash\psi\) or \(T(A)\vdash\neg\psi\). Finally \(T\subseteq T(A)\)
by definition.

\emph{Clause 2.} The \(L\)-reduct \(M_L\) of \(M\) is a model of \(T\),
because \(T\subseteq T'\) and \(T\) is a theory of \(L\). So \(A\), its
subsystem on \(\{c^M:c\in C\}\), belongs to \(\mathcal N(T)\) (when
\(C\) is empty, \(A\) is the empty system, which the standing convention
of \cref{def:p03-morley-standing} admits, and \(T(A)=T\)). By the
argument of clause~1, \(N=(M_L,a)_{a\in|A|}\) is a model of \(T(A)\). In
\(M\) the constant \(c\) denotes \(c^M\), and in \(N\) the constant
\(\overline{c^M}\) denotes the same element, while \(M\) and \(N\) agree
on \(L\); by induction on formulas, \(M\entails\varphi\) if and only if
\(N\entails\sigma\varphi\) for every sentence \(\varphi\) of \(L\cup C\).
Because \(T'\) is complete and \(M\entails T'\), \(T'\vdash\varphi\) if
and only if \(M\entails\varphi\): if \(T'\) does not prove \(\varphi\) it
proves \(\neg\varphi\). In the same way, because \(T(A)\) is complete by
clause~1 and \(N\entails T(A)\), \(T(A)\vdash\sigma\varphi\) if and only if
\(N\entails\sigma\varphi\). Combining the three equivalences gives the
claim.

For a second model \(M'\) of \(T'\), and \(c,d\in C\),
\(c^M=d^M\) if and only if \(T'\vdash c=d\), if and only if
\(c^{M'}=d^{M'}\), by the completeness of \(T'\) as above. So
\(h(c^M)=c^{M'}\) defines a one-one map of \(|A|\) onto \(|A'|\). For a
relation symbol \(R\) of degree \(k\) and \(c_1,\dots,c_k\in C\),
\(R^A(c_1^M,\dots,c_k^M)\) holds if and only if
\(T'\vdash R(c_1,\dots,c_k)\), if and only if
\(R^{A'}(c_1^{M'},\dots,c_k^{M'})\) holds, because \(A\) and \(A'\) are
induced subsystems. Every tuple of \(|A|\) has this form, so \(h\)
preserves and reflects every relation and is an isomorphism.

\emph{Clause 3.} Suppose \(f\) is a monomorphism and
\(\psi\in\mathscr D(A)\). An atomic sentence \(R(\bar a_1,\dots,\bar a_k)\)
is true in \((A,a)_{a\in|A|}\) if and only if
\(R(\overline{f(a_1)},\dots,\overline{f(a_k)})\) is true in
\((B,b)_{b\in|B|}\), because \(f\) preserves and reflects \(R\); and
\(\bar a=\bar a'\) is true if and only if \(f(a)=f(a')\), because \(f\)
is one-one. By induction on Boolean combinations, \(\tilde f\psi\) is a
true open sentence of \(L(B)\), so \(\tilde f\psi\in\mathscr D(B)\subseteq
T(B)\).

Suppose next that \(T(B)\vdash\tilde f\psi\) for every
\(\psi\in\mathscr D(A)\), let \(N\) be a model of \(T(B)\), and let
\(N_f\) be the \(L(A)\)-structure with the \(L\)-reduct of \(N\) in which
\(\bar a\) denotes the element that \(\overline{f(a)}\) denotes in \(N\).
For every sentence \(\theta\) of \(L(A)\), \(N_f\entails\theta\) if and
only if \(N\entails\tilde f\theta\). The sentences of \(T\) contain no
constants of \(L(A)\setminus L\), so \(\tilde f\) fixes them and
\(N_f\entails T\); and \(N\entails\tilde f\psi\) for
\(\psi\in\mathscr D(A)\), so \(N_f\entails\mathscr D(A)\). Thus
\(N_f\entails T(A)\). If \(T(A)\vdash\theta\), then \(N_f\entails\theta\)
and \(N\entails\tilde f\theta\). As \(N\) was an arbitrary model of
\(T(B)\), \(T(B)\vdash\tilde f\theta\) by the completeness theorem
(\cref{thm:fol-completeness}).

Suppose finally that \(T(B)\vdash\tilde f\theta\) whenever
\(T(A)\vdash\theta\). In particular this holds for every
\(\theta\in\mathscr D(A)\subseteq T(A)\). Let \(M\) be a model of \(T\)
with \(B\subseteq M\); by the argument of clause~1,
\(N=(M,b)_{b\in|B|}\) is a model of \(T(B)\). If \(R^A(a_1,\dots,a_k)\)
holds, then \(R(\bar a_1,\dots,\bar a_k)\in\mathscr D(A)\), so
\[
 N\entails R(\overline{f(a_1)},\dots,\overline{f(a_k)});
\]
thus
\(R^M(f(a_1),\dots,f(a_k))\) holds, and so does
\(R^B(f(a_1),\dots,f(a_k))\), because \(B\) is a subsystem of \(M\). If
\(R^A(a_1,\dots,a_k)\) fails, the negated atomic sentence is in
\(\mathscr D(A)\), and the same argument shows that
\(R^B(f(a_1),\dots,f(a_k))\) fails. If \(a\neq a'\), then
\(\neg\,\bar a=\bar a'\) is in \(\mathscr D(A)\), so
\(f(a)\neq f(a')\). Hence \(f\) is a monomorphism.

\emph{Clause 4.} A point \(p\in S(A)\) is a maximal consistent set of
formulas of \(F(A)\), closed under equivalence in \(T(A)\); for each
\(\psi\in F(A)\) exactly one of \(\psi\) and \(\neg\psi\) is in \(p\).
Put \(T_p=T(A)\cup p(c)\).

\(T_p\) is consistent: by \cref{thm:p03-morley-types} and the remark
opening this annex, some model \(M\supseteq A\) of \(T\) has an element
\(b\) realizing \(p\), so every formula of \(p\) holds of \(b\) in
\((M,a)_{a\in|A|}\), and \((M,a,b)_{a\in|A|}\), with \(c\) denoting \(b\),
is a model of \(T_p\). \(T_p\) is complete: let \(\theta\) be a sentence
of \(L(A)\cup\{c\}\), rename its bound variables so that \(v_0\) does not
occur bound, and let \(\psi\) be the result of replacing \(c\) by
\(v_0\). Then \(\psi\in F(A)\), and \(\psi(c)\) is \(\theta\) up to the
renaming of bound variables. Since \(\psi\in p\) or \(\neg\psi\in p\),
\(T_p\vdash\theta\) or \(T_p\vdash\neg\theta\).

The map is one-one: if \(p\neq p'\), some \(\psi\) lies in \(p\) and not
in \(p'\), so \(\neg\psi\in p'\); then \(T_p\vdash\psi(c)\) and
\(T_{p'}\vdash\neg\psi(c)\), and since \(T_{p'}\) is consistent the two
theories have different consequences.

The map is onto: let \(T''\supseteq T(A)\) be a consistent complete theory
of \(L(A)\cup\{c\}\), and put \(p=\{\psi\in F(A):T''\vdash\psi(c)\}\). If
\(T(A)\vdash\forall v_0(\psi\to\psi')\) and \(\psi\in p\), then
\(\psi'\in p\); \(p\) is closed under conjunction; and, \(T''\) being
consistent and complete, \(p\) contains exactly one of \(\psi\) and
\(\neg\psi\) for each \(\psi\in F(A)\). So the classes of the members of
\(p\) form an ultrafilter of \(F(A)\), that is, \(p\in S(A)\). Every
member of \(T_p\) is a consequence of \(T''\). If \(T''\vdash\theta\),
then \(T_p\) cannot prove \(\neg\theta\), since \(T''\) would then be
inconsistent; as \(T_p\) is complete, \(T_p\vdash\theta\). So \(T_p\) and
\(T''\) have the same consequences.
\end{proof}
\end{annexitem}

\begin{annexitem}{annex:p03-w03-p03-033-open-formulas}{thm:p03-w03-p03-033-open-formulas}{Open formulas and satisfaction}
\begin{proof}
\emph{Clause 1.} Let \(\psi\in F(A)\), let \(\bar a_1,\dots,\bar a_n\) be
the distinct constants of \(L(A)\setminus L\) in \(\psi\), and let
\(\chi(v_0,w_1,\dots,w_n)\) be the \(L\)-formula obtained by replacing
them by distinct variables \(w_1,\dots,w_n\) not occurring in \(\psi\).
The free variables of \(\chi\) are among \(v_0,w_1,\dots,w_n\), so
Lemma~1.1(e) (\cref{thm:p03-morley-normalization}) gives a relation
symbol \(R\) of degree \(n+1\) with
\(T\vdash\forall v_0\forall w_1\cdots\forall w_n
(\chi\leftrightarrow R(v_0,w_1,\dots,w_n))\). Instantiating \(w_i\) by
\(\bar a_i\) gives
\[
 T(A)\vdash\forall v_0\bigl(\psi\leftrightarrow R(v_0,\bar a_1,\dots,\bar a_n)\bigr),
\]
and \(R(v_0,\bar a_1,\dots,\bar a_n)\) is an open formula of \(F(A)\).

Let \(F_o(A)\) be the set of open formulas in \(F(A)\). It is closed under
\(\wedge\), \(\vee\) and \(\neg\), and two of its members are identified
exactly when they are identified in \(F(A)\), since the identification is
in both cases equivalence in \(T(A)\). So the identification turns
\(F_o(A)\) into a Boolean subalgebra of \(F(A)\), and the inclusion is
one-one. By the first paragraph every class of \(F(A)\) contains a member
of \(F_o(A)\), so the inclusion is onto and is an isomorphism of Boolean
algebras. An isomorphism of Boolean algebras carries ultrafilters to
ultrafilters and basic clopen sets to basic clopen sets, so
\(p\mapsto p\cap F_o(A)\) is a homeomorphism of \(S(A)\) onto the Stone
space of \(F_o(A)\).

\emph{Clause 2.} If \(A\entails\psi\), then \(\psi\in\mathscr D(A)\subseteq
T(A)\). If \(A\not\entails\psi\), then \(\neg\psi\) is an open sentence true
in \((A,a)_{a\in|A|}\), so \(\neg\psi\in\mathscr D(A)\) and
\(T(A)\vdash\neg\psi\); since \(T(A)\) is consistent
(\cref{thm:p03-w03-p03-033-syntactic}), \(T(A)\) does not prove \(\psi\).

\emph{Clause 3.} If \(A\) is a model of \(T\), then \((A,a)_{a\in|A|}\) is
a model of \(T\cup\mathscr D(A)=T(A)\). Hence \(T(A)\vdash\psi\) implies
\(A\entails\psi\); and if \(T(A)\) does not prove \(\psi\), then, being
complete (\cref{thm:p03-w03-p03-033-syntactic}), it proves \(\neg\psi\), so
\(A\not\entails\psi\).

\emph{Clause 4.} If \(A\) is not a model of \(T\), some \(\sigma\in T\) is
false in \(A\). It is a sentence of \(L\), and \(T(A)\vdash\sigma\) since
\(\sigma\in T(A)\).

\emph{The example.} Let \(M\) be a model of \(T\) with \(A\subseteq M\),
so \(a<^Mb\). The theory \(\Sigma=\mathsf{DLO}\) contains density, and
\(T=\Sigma^*\) contains every sentence of the original language true in
all models of \(\Sigma\); so \(T\vdash\forall v_0\forall v_1(v_0<v_1\to
\exists v_2(v_0<v_2\wedge v_2<v_1))\). With \(\bar a<\bar b\in\mathscr
D(A)\) this gives \(T(A)\vdash\psi\), after renaming the bound variable.
In \(A\), whose universe is \(\{a,b\}\), the element \(a\) is not above
\(a\) and \(b\) is not below \(b\), because \(<\) is irreflexive in
\(M\) and \(A\) is a subsystem; so no element of \(|A|\) lies strictly
between \(a\) and \(b\), and \(A\not\entails\psi\). The relation \(R\) is
interpreted in \(M\) by the truth set of
\(\exists v_2(v_0<v_2\wedge v_2<v_1)\), which contains \((a,b)\) by
density; as \(A\) is a subsystem of \(M\), \(R^A(a,b)\) holds, so
\(R(\bar a,\bar b)\in\mathscr D(A)\). \(A\) is not a model of \(T\), since
it fails the density sentence of \(\Sigma\).
\end{proof}
\end{annexitem}

\begin{annexitem}{annex:p03-w03-p03-033-tarski-criterion}{thm:p03-w03-p03-033-tarski-criterion}{Tarski's criterion from the Tarski--Vaught test}
\begin{proof}
By Morley's definition, \(i_{AB}\) is an elementary monomorphism when every
formula of \(L\) true of a tuple from \(|A|\) in \(A\) is true of the
same tuple in \(B\); applying this to negations as well, it says that
\((A,a)_{a\in|A|}\equiv(B,a)_{a\in|A|}\), that is, \(A\preccurlyeq B\) as
structures of the relational language \(L\). Since \(A\subseteq B\), \(A\)
is a substructure of \(B\). By the Tarski--Vaught test
(\cref{thm:tarski-vaught-fol}), \(A\preccurlyeq B\) if and only if for
every \(L\)-formula \(\varphi(y,\bar x)\) and tuple \(\bar a\) from
\(|A|\), \(B\entails\exists y\,\varphi(y,\bar a)\) implies
\(B\entails\varphi(b,\bar a)\) for some \(b\in|A|\). Renaming \(y\) as
\(v_0\) and the variables \(\bar x\) so that \(v_0\) is the only other free
variable, the formulas \(\varphi(v_0,\bar a)\) are exactly the members of
\(F(A)\): each member of \(F(A)\) has this form, with \(\varphi\) obtained
by replacing its constants by fresh variables. The hypothesis
\(B\entails\exists v_0\,\varphi(v_0,\bar a)\) says that some element of
\(|B|\) satisfies \(\varphi(v_0,\bar a)\) in \(B\), and the conclusion
says that some element of \(|A|\) does. So the two conditions coincide.

If \(i_{AB}\) is elementary, every \(\sigma\in T\) is true in \(B\) and
hence in \(A\), so \(A\) is a model of \(T\).
\end{proof}
\end{annexitem}

\begin{annexitem}{annex:p03-w03-p03-033-variant-rank}{thm:p03-w03-p03-033-variant-rank}{Monotonicity and existence of the single-point rank}
\begin{proof}
\emph{A fibre lemma.} Let \(f:A\to B\), let
\(p\in\operatorname{Tr}_\natural^\gamma(A)\), and let
\(q\in(f^*)^{-1}(p)\cap S_\natural^\gamma(B)\). Then
\(q\in\operatorname{Tr}_\natural^\gamma(B)\). Indeed, for a monomorphism
\(g:B\to C\), if \(g^*(r)=q\) then \((g\circ f)^*(r)=f^*(q)=p\), so
\[
 (g^*)^{-1}(q)\cap S_\natural^\gamma(C)\subseteq
 ((g\circ f)^*)^{-1}(p)\cap S_\natural^\gamma(C),
\]
and the right-hand set has at most one point because
\(p\in\operatorname{Tr}_\natural^\gamma(A)\) and \(g\circ f\) is a
monomorphism.

\emph{Clause 1.} Let \(\rho_\natural(p)=\beta\) and
\(q\in(f^*)^{-1}(p)\). If \(q\in S_\natural^{\beta+1}(B)\), then
\(q\in S_\natural^\beta(B)\), and the fibre lemma with \(\gamma=\beta\)
puts \(q\) in \(\operatorname{Tr}_\natural^\beta(B)\), which is disjoint
from \(S_\natural^{\beta+1}(B)\). So \(q\notin S_\natural^{\beta+1}(B)\),
that is, \(q\in\operatorname{Tr}_\natural^\gamma(B)\) for some
\(\gamma\leq\beta\), and \(\rho_\natural(q)\leq\beta\). A fibre point of
single-point rank exactly \(\beta\) lies in
\((f^*)^{-1}(p)\cap S_\natural^\beta(B)\), which has at most one point by
the definition of \(\operatorname{Tr}_\natural^\beta(A)\).

\emph{Clause 2.} Let \(q\in S_\natural^\beta(B)\) and \(p=f^*(q)\).
If \(p\notin S_\natural^\beta(A)\), then
\(p\in\operatorname{Tr}_\natural^\gamma(A)\) for some \(\gamma<\beta\).
Since \(q\in S_\natural^\beta(B)\subseteq S_\natural^\gamma(B)\), the fibre
lemma gives \(q\in\operatorname{Tr}_\natural^\gamma(B)\), contradicting
\(q\in S_\natural^\beta(B)\).

\emph{Isomorphism invariance.} Let \(h:B\to B'\) be a monomorphism onto
\(B'\); its inverse is a monomorphism, and \(h^*\) and \((h^{-1})^*\) are
mutually inverse bijections by functoriality. Clause~2 applied to \(h\) and
to \(h^{-1}\) shows that \(h^*\) maps \(S_\natural^\beta(B')\) onto
\(S_\natural^\beta(B)\) for every \(\beta\). Since
\(\operatorname{Tr}_\natural^\beta=S_\natural^\beta\setminus
S_\natural^{\beta+1}\), it also maps \(\operatorname{Tr}_\natural^\beta(B')\)
onto \(\operatorname{Tr}_\natural^\beta(B)\). Hence
\(\rho_\natural(h^*(q'))=\rho_\natural(q')\), with both sides undefined
together.

\emph{A set of values.} Fix an infinite cardinal \(\lambda\). Every
\(B\in\mathcal N(T)\) with \(\kappa(B)\leq\lambda\) is isomorphic to a
system whose universe is a subset of \(\lambda\), and these systems form a
set, since the language is countable. By invariance, the single-point
ranks of points over systems of size at most \(\lambda\) form a set
\(V_\lambda\) of ordinals. Let \(\beta_\lambda\) be the least ordinal
greater than every member of \(V_\lambda\).

\emph{Small witnesses.} Suppose \(p\in S_\natural^\beta(A)\setminus
\operatorname{Tr}_\natural^\beta(A)\). Then some \(f:A\to B\) has distinct
points \(q_1,q_2\in(f^*)^{-1}(p)\cap S_\natural^\beta(B)\). Choose a formula
\(\varphi\in q_1\setminus q_2\), let \(E\subseteq B\) be its finite set of
parameters, and let \(B_0\) be the induced subsystem on \(f(A)\cup E\),
with inclusion \(i:B_0\to B\) and \(f_0:A\to B_0\) the map \(f\) with
restricted range. Then \(\varphi\) belongs to \(i^*(q_1)\) and not to
\(i^*(q_2)\), so these points differ. Both lie in \(S_\natural^\beta(B_0)\)
by clause~2, both restrict to \(f_0^*i^*(q_j)=f^*(q_j)=p\), and
\(\kappa(B_0)\leq\kappa(A)+\aleph_0\).

\emph{Clause 3.} Let \(T\) be totally transcendental. Every point \(q\)
then has an ordinal rank \(\alpha_q\) and a degree \(d(q)\geq1\)
(\cref{def:p03-morley-total-transcendence,thm:p03-morley-degree}). The
pairs \((\alpha_q,d(q))\), ordered lexicographically, are well ordered.
Suppose some point has no single-point rank, and choose such a point
\(p\in S(A)\) whose pair \((\alpha,d)\) is least among all such points over
all systems. Put \(\lambda=\kappa(A)+\aleph_0\) and \(\beta=\beta_\lambda\).
Then \(p\in S_\natural^\beta(A)\setminus\operatorname{Tr}_\natural^\beta(A)\),
so by small witnesses there are \(f_0:A\to B_0\) with
\(\kappa(B_0)\leq\lambda\) and distinct
\(q_1,q_2\in(f_0^*)^{-1}(p)\cap S_\natural^\beta(B_0)\). If \(q_j\) had a
single-point rank, it would belong to \(V_\lambda\) and so be less than
\(\beta\); but membership in \(S_\natural^\beta(B_0)\) excludes every
\(\operatorname{Tr}_\natural^\gamma(B_0)\) with \(\gamma<\beta\). So neither
\(q_1\) nor \(q_2\) has a single-point rank.

Compare their pairs with \((\alpha,d)\). Since \(q_j\in
S^{\alpha_{q_j}}(B_0)\) and \(f_0^*(S^{\alpha_{q_j}}(B_0))=
S^{\alpha_{q_j}}(A)\) by \cref{thm:p03-morley-derivative}, the point \(p\)
lies in \(S^{\alpha_{q_j}}(A)\); as \(p\in\operatorname{Tr}^\alpha(A)\),
this forces \(\alpha_{q_j}\leq\alpha\). If both ranks equal \(\alpha\), both
points lie in \((f_0^*)^{-1}(p)\cap S^\alpha(B_0)\), whose degrees add up
to \(d\) by \cref{thm:p03-morley-degree}; then
\(d(q_1)+d(q_2)\leq d\) and each \(d(q_j)<d\). Otherwise one of the ranks
is smaller than \(\alpha\). In both cases one of \(q_1,q_2\) is a point
without single-point rank whose pair is lexicographically smaller than
\((\alpha,d)\), contradicting the choice of \(p\).
\end{proof}
\end{annexitem}

\begin{annexitem}{annex:p03-w03-p03-033-variant-field}{thm:p03-w03-p03-033-variant-field}{Two algebraic types}
\begin{proof}
\emph{Uniquely solvable formulas.} Let \(q\in S(B)\) contain a formula
\(\psi(v_0)\) over \(B\) such that
\(\forall v_0\forall v_1\,(\psi(v_0)\wedge\psi(v_1)\to v_0=v_1)\) holds in
some model containing \(B\). That sentence over \(B\) is equivalent to an
atomic relation of its parameters, so it holds in every model into which
\(B\) is mapped by a monomorphism. Let \(g:B\to C\) and let \(r_1,r_2\) be
points of \((g^*)^{-1}(q)\). Realize them by \(e_1,e_2\) in one model
containing \(C\). Both satisfy the translate of \(\psi\), so \(e_1=e_2\)
and \(r_1=r_2\). Thus every fibre has at most one point, and
\(q\in\operatorname{Tr}_\natural^0(B)\), because
\(S_\natural^0(B)=S(B)\).

For \(B\in\mathcal N(T)\) let \(K(B)\) be the subfield generated by
\(B\) in a model containing it. An element \(c\in K(B)\) is a quotient
\(r(\bar b)/s(\bar b)\) of polynomials over \(\mathbb Z\) evaluated at
parameters, with \(s(\bar b)\ne0\). The formula
\(s(\bar b)\,v_0=r(\bar b)\) is uniquely solvable in every field, so the
type of \(c\) over \(B\) lies in \(\operatorname{Tr}_\natural^0(B)\).

\emph{The types are not of single-point rank zero.} For \(k=2,3\) let
\(r\) and \(r'\) be distinct roots of \(v^k-2\) in an algebraically closed
field of characteristic \(0\), and let \(B\) be the induced system on
\(\{0,r\}\). The types of \(r\) and \(r'\) over \(B\) both contain
\(v_0^k=2\), so both restrict to \(p_k\), the only type over \(A\)
containing that equation (\cref{thm:p03-morley-field-example}). They
differ, since only the first contains \(v_0=r\). So the fibre of \(p_k\)
over \(B\) has two points of \(S_\natural^0(B)=S(B)\), and
\(p_k\notin\operatorname{Tr}_\natural^0(A)\).

\emph{The types have single-point rank one.} Let \(f:A\to B\) be any
monomorphism. Realize all points of \((f^*)^{-1}(p_k)\) in one model
\(D\supseteq B\). Each realizer satisfies \(v_0^k=2\), so the fibre points
are the types over \(B\) of the roots of \(v^k-2\) in \(D\). Put
\(K=K(B)\subseteq D\). Roots in \(K\) have types in
\(\operatorname{Tr}_\natural^0(B)\), by the previous paragraph. The roots
outside \(K\) are the roots of a single polynomial irreducible over \(K\).
For \(k=2\) the roots are \(\pm c\), so either both lie in \(K\) or
\(v^2-2\) has no root in \(K\) and, being of degree \(2\), is irreducible
over \(K\). For \(k=3\) the three roots add up to \(0\), the coefficient
of \(v^2\) being \(0\); so if two roots lie in \(K\) the third does, and
the number of roots in \(K\) is \(0\), \(1\) or \(3\). With no root in
\(K\), the cubic is irreducible over \(K\), because a reducible cubic has
a linear factor. With exactly one root \(c\in K\), the quotient
\((v^3-2)/(v-c)\) is a quadratic over \(K\) with no root in \(K\), hence
irreducible. With three roots in \(K\) nothing remains outside. An
irreducible polynomial over \(K(B)\), written over \(B\) after clearing
denominators, is contained in exactly one type over \(B\)
(\cref{thm:p03-morley-field-example}), so all roots outside \(K\) have the
same type. Consequently
\((f^*)^{-1}(p_k)\cap S_\natural^1(B)\) has at most one point. Together
with the previous step, \(p_k\in S_\natural^1(A)\) and
\(p_k\in\operatorname{Tr}_\natural^1(A)\), so
\(\rho_\natural(p_2)=\rho_\natural(p_3)=1\), while \(d(p_2)=2\) and
\(d(p_3)=3\).

\emph{The reading with exactly one point.} Define
\(S_=^\beta\) and \(\operatorname{Tr}_=^\beta\) by the same recursion, with
``at most one point'' replaced by ``exactly one point''. Restriction maps
are onto (\cref{thm:p03-morley-types}), so a fibre in \(S(B)\) is never
empty. Hence \(\operatorname{Tr}_=^0=\operatorname{Tr}_\natural^0\) and
\(S_=^1=S_\natural^1\) for every system. Let \(r_1,r_2,r_3\) be the three
roots of \(v^3-2\) in an algebraically closed field of characteristic
\(0\), and let \(B_3\) be the induced system on \(\{0,r_1,r_2,r_3\}\). A
point of the fibre of \(p_3\) over \(B_3\) is realized, in a model
containing \(B_3\), by a root of \(v^3-2\), and that model's only such
roots are \(r_1,r_2,r_3\). So every fibre point contains some
\(v_0=r_i\) and lies in \(\operatorname{Tr}_\natural^0(B_3)=
\operatorname{Tr}_=^0(B_3)\). Therefore the fibre meets \(S_=^1(B_3)\) in
no point, and, the stages being decreasing, it meets no
\(S_=^\beta(B_3)\) with \(\beta\geq1\). For \(\beta\geq1\) the inclusion
of \(A\) in \(B_3\) thus violates the condition for
\(p_3\in\operatorname{Tr}_=^\beta(A)\). The extension over \(\{0,r\}\)
above violates it for \(\beta=0\). So \(p_3\) lies in every \(S_=^\beta(A)\)
and receives no rank. The theory \(\mathrm{ACF}_0\) is totally
transcendental, since its second derivative is empty
(\cref{thm:p03-morley-field-example}).
\end{proof}
\end{annexitem}

\begin{annexitem}{annex:p03-w03-p03-033-automorphism-test}{thm:p03-w03-p03-033-automorphism-test}{Automorphisms by back and forth}
\begin{proof}
\emph{Sufficiency of the weakened condition, hence clause~1.} Let
\(p,p'\in S(A)\) both contain \(Q\). Realize them by \(b\) and \(b'\) in
one model \(B\supseteq A\). Both satisfy \(Q\), so there are a model
\(C\supseteq B\) and an automorphism \(\sigma\) of \(C\) fixing \(A\)
pointwise with \(\sigma(b)=b'\). Since \(B\) is an induced subsystem of
\(C\), \(b\) and \(b'\) still realize \(p\) and \(p'\) in \(C\). Every
\(\varphi\in F(A)\) is equivalent to an atomic relation \(R(v_0,\bar a)\)
with \(\bar a\) from \(A\), and \(\sigma\) preserves and reflects
\(R\) while fixing \(\bar a\). So \(\varphi\in p\) if and only if
\(\varphi\in p'\), and \(p=p'\). Clause~1 is the case \(C=B\).

\emph{An extension lemma.} Let \(C\) be a model, \(D\subseteq C\) an
induced subsystem and \(h:D\to C\) a monomorphism. Then there are a model
\(C'\supseteq C\) and a monomorphism \(u:C\to C'\) extending \(h\).
By amalgamation over a common subsystem
(\cref{thm:p03-morley-normalization}), applied to the inclusion
\(D\to C\) and to \(h:D\to C\), there are \(E\in\mathcal N(T)\) and
monomorphisms \(u,v:C\to E\) with \(u(x)=v(h(x))\) for \(x\in D\).
Replacing \(E\) by an isomorphic copy, we may take \(v\) to be an
inclusion \(C\subseteq E\); then \(u\) extends \(h\). Finally \(E\) is a
subsystem of some model \(C'\).

\emph{Necessity of the weakened condition.} Suppose \(Q\) lies in exactly
one point \(p\). Let \(B\supseteq A\) be a model and let \(b,b'\in B\)
satisfy \(Q\). The types of \(b\) and \(b'\) over \(A\) contain \(Q\), so
both equal \(p\), and by \cref{thm:p03-morley-types} the map \(h_0\) that
is the identity on \(A\) and sends \(b\) to \(b'\) is a monomorphism of
the induced system on \(A\cup\{b\}\) into \(B\). Put \(C_0=B\).
Apply the extension lemma to \(h_0\) to get a model \(C_1\supseteq C_0\)
and \(F_0:C_0\to C_1\) extending \(h_0\). Recursively, given
\(F_k:C_{2k}\to C_{2k+1}\), apply the lemma to the monomorphism
\(F_k^{-1}\), defined on the induced subsystem \(F_k(C_{2k})\) of
\(C_{2k+1}\) with values in \(C_{2k}\subseteq C_{2k+1}\), to get a model
\(C_{2k+2}\supseteq C_{2k+1}\) and \(G_k:C_{2k+1}\to C_{2k+2}\) extending
\(F_k^{-1}\). Then apply it to \(G_k^{-1}\), defined on
\(G_k(C_{2k+1})\subseteq C_{2k+2}\), to get
\(C_{2k+3}\supseteq C_{2k+2}\) and \(F_{k+1}:C_{2k+2}\to C_{2k+3}\)
extending \(G_k^{-1}\).

Since \(G_k\) extends \(F_k^{-1}\), it maps \(F_k(C_{2k})\) onto
\(C_{2k}\), so \(G_k^{-1}\) restricted to \(C_{2k}\) is \(F_k\), and
\(F_{k+1}\) extends \(F_k\). The range of \(F_{k+1}\) contains the range
\(C_{2k+1}\) of \(G_k^{-1}\). Let \(C=\bigcup_nC_n\), a model by
\cref{thm:p03-morley-normalization}, and \(F=\bigcup_kF_k\). Then \(F\) is
defined on \(\bigcup_kC_{2k}=C\), is a monomorphism as a union of a chain
of monomorphisms, and is onto \(C\) because its range contains every
\(C_{2k+1}\). So \(F\) is an automorphism of \(C\); it extends \(h_0\),
fixes \(A\) pointwise and sends \(b\) to \(b'\), and \(C\supseteq B\).
\end{proof}
\end{annexitem}

\begin{annexitem}{annex:p03-w03-p03-033-dlo-density}{thm:p03-w03-p03-033-dlo-density}{Isolated cut types in dense linear orders}
\begin{proof}
\emph{Types are order positions.} By \cref{thm:p03-w03-p03-022-dlo-step},
\(\mathsf{DLO}\) eliminates quantifiers. A formula of \(T\) over \(A\) in
the variable \(v_0\) is equivalent modulo \(T\) to a \(\mathsf{DLO}\)
formula \(\varphi(v_0,\bar a)\), hence modulo \(\mathsf{DLO}\) to a
quantifier-free \(\psi(v_0,\bar a)\) whose atomic parts compare \(v_0\)
and the finitely many \(a_i\) by \(<\) and \(=\). Consequently, if \(c\)
and \(c'\), in models containing \(A\), stand in the same relations
\(<,=,>\) to every element of \(A\), they satisfy the same formulas over
\(A\) and have the same type.

\emph{Density.} Let \(U_\varphi\subseteq S(A)\) be a nonempty basic
open set and choose \(p\in U_\varphi\), realized by \(b\) in a model
\(M\supseteq A\). If \(b\in A\), then \(p\) contains \(v_0=b\), whose only
solution is \(b\); so \(p\) is isolated. Otherwise let
\(a_1<\dots<a_k\) be the parameters of the quantifier-free form \(\psi\)
of \(\varphi\), and let \(I\) be the open interval of \(M\) between the
largest \(a_i\) below \(b\) and the smallest \(a_i\) above \(b\), each end
omitted when there is no such \(a_i\); for \(k=0\), \(I\) is all of \(M\).

If some \(a\in A\) lies in \(I\), then \(a\) differs from every \(a_i\)
and has the same order relations to them as \(b\). Since \(\psi\) only
mentions those parameters, \(\psi(a,\bar a)\) holds, and the type of \(a\)
over \(A\) is an isolated point of \(U_\varphi\).

If no element of \(A\) lies in \(I\), let \(\chi(v_0)\) be the formula
saying that \(v_0\) lies strictly between the two ends of \(I\) that are
present (the formula \(v_0=v_0\) when neither is). Then \(p\) contains
\(\chi\). Any \(c\) satisfying \(\chi\), in a model containing \(A\), has
the same relations to \(A\) as \(b\): an element \(a\in A\) is not in
\(I\), so it is at most the lower end, hence below both \(b\) and \(c\),
or at least the upper end, hence above both. By the first paragraph
every point of \(U_\chi\) equals \(p\); so \(U_\chi=\{p\}\) and \(p\) is an
isolated point of \(U_\varphi\). This covers the empty system, where
\(k=0\) and no \(a\) exists.

\emph{Not totally transcendental.} By \cref{thm:p03-morley-thin-models}
the countable rational model has continuum many types over itself. A
totally transcendental theory has
\(\kappa(S(A))\leq\kappa(A)+\aleph_0\) by
\cref{thm:p03-morley-countability}, which is countable here. So \(T\) is
not totally transcendental.

\emph{Consequences.} The density shown above is the hypothesis of the
prime-model clause of \cref{thm:p03-morley-prime-models}, so every
subsystem has a prime model over it. The chain theorem
\cref{thm:p03-morley-chain-primes} is stated for totally transcendental
theories, and that hypothesis fails.
\end{proof}
\end{annexitem}

\begin{annexitem}{annex:p03-w03-p03-033-endomorphism}{thm:p03-w03-p03-033-endomorphism}{An indiscernible set, compactness, and a prime model}
\begin{proof}
\emph{An infinite indiscernible set.} The language of \(T\) is countable
and \(T\) has an infinite model, so \(T\) has a model \(B\) of power
\(\aleph_1\) (\cref{thm:p03-morley-thin-models}). Choose \(a\in B\) and let
\(A_0\) be the induced system on \(\{a\}\). Since
\(\kappa(A_0)<\kappa(B)=\aleph_1\) and \(\aleph_1\) is regular and
uncountable, \cref{thm:p03-morley-indiscernible-extraction} gives
\(X_0\subseteq|B|\setminus\{a\}\) with \(\kappa(X_0)=\aleph_1\),
indiscernible over \(A_0\). If \(\tau:X_0\to X_0\) is one-one, extending
it by the identity on \(\{a\}\) gives a one-one map of
\(\{a\}\cup X_0\) into itself fixing \(A_0\), hence a monomorphism; its
restriction \(\tau\) is a monomorphism of the induced system on \(X_0\).

\emph{Stretching to power \(\kappa\).} Add new constants \(c_\xi\) for
\(\xi<\kappa\). Let \(\Delta\) consist of \(T\), the sentences
\(c_\xi\ne c_\eta\) for \(\xi\ne\eta\), and, for every relation symbol
\(R\) of \(T\) of arity \(n\) and all index tuples
\((\xi_1,\dots,\xi_n)\) and \((\eta_1,\dots,\eta_n)\) with the same
equality pattern (\(\xi_i=\xi_j\) exactly when \(\eta_i=\eta_j\)), the
sentence
\(R(c_{\xi_1},\dots,c_{\xi_n})\leftrightarrow
  R(c_{\eta_1},\dots,c_{\eta_n})\).
A finite \(\Delta_0\subseteq\Delta\) mentions constants with indices in a
finite set \(J\). Interpret them in \(B\) by a one-one \(e:J\to X_0\). For
a sentence of the third kind, the assignment \(e(\xi_i)\mapsto e(\eta_i)\)
is well defined and one-one by the equality condition; it maps a finite
subset of \(X_0\) into \(X_0\), and it extends to a one-one map of \(X_0\)
into itself, because the two complements in \(X_0\) of the finite domain
and range both have power \(\aleph_1\). That extension is a monomorphism
of \(X_0\), so the sentence holds in \(B\). Hence \(\Delta_0\) has a model,
and by compactness (\cref{thm:p03-compactness-ultraproduct}) so has
\(\Delta\). Let \(C\) be a model of \(\Delta\), \(C_0\) its reduct, a
model of \(T\), and \(X=\{c_\xi^C:\xi<\kappa\}\), so \(\kappa(X)=\kappa\).
For a one-one \(\tau:X\to X\) and a tuple \(\bar x=c_{\bar\xi}^C\) from
\(X\), the tuple \(\tau(\bar x)\) is \(c_{\bar\eta}^C\) with the same
equality pattern, so \(\Delta\) gives
\(R(\bar x)\leftrightarrow R(\tau(\bar x))\) in \(C_0\). Thus every
one-one map of \(X\) into itself is a monomorphism of the induced system
\(X\in\mathcal N(T)\).

\emph{The prime model.} Since \(T\) is totally transcendental,
\cref{thm:p03-morley-prime-models} gives a model \(A\supseteq X\) prime
over \(X\) with \(\kappa(A)\leq\kappa(X)+\aleph_0=\kappa\); and
\(\kappa(A)\geq\kappa(X)\), so \(\kappa(A)=\kappa\). Let \(\tau:X\to X\)
be one-one. Followed by the inclusion of \(X\) in \(A\), it is a
monomorphism of \(X\) into the model \(A\). By primeness it extends to a
monomorphism \(g:A\to A\), an endomorphism of \(A\) whose restriction to
\(X\) is \(\tau\).
\end{proof}
\end{annexitem}

\begin{annexitem}{annex:p03-w03-p03-033-question-five}{thm:p03-w03-p03-033-question-five}{Total indiscernibility from the exclusion of order}
\begin{proof}
\emph{Order indiscernibles.} The theory \(\Sigma\) is complete with an
infinite model, so all models of \(T\) are infinite; models of \(T'\) are
models of \(T^\#\), whose reducts are models of \(T\), so \(T'\) has an
infinite model. Let \(X\) be the ordinal \(\kappa\) with its order. By
\cref{thm:p03-morley-ramsey-em} there is a model \(D\) of \(T'\) with
\(X\subseteq|D|\) a set of order indiscernibles for all
\(L^\#\)-formulas.

\emph{Total indiscernibility.} Let \(\psi(w_0,\dots,w_{m-1})\) be an
\(L^\#\)-formula. For \(m\leq1\), any two tuples of \(m\) distinct
elements of \(X\) are increasing, so they agree on \(\psi\). Let
\(m\geq2\). For a permutation \(s\) of \(\{0,\dots,m-1\}\) and an
increasing tuple \(\bar x\) from \(X\), write
\(\bar x^s=(x_{s(0)},\dots,x_{s(m-1)})\). Applying order indiscernibility
to the formula \(\psi(w_{s(0)},\dots,w_{s(m-1)})\) shows that the set
\(P=\{s:D\models\psi(\bar x^s)\}\) does not depend on \(\bar x\). If \(P\)
were neither empty nor the whole symmetric group, the relation defined by
\(\psi\) would be connected and antisymmetric on \(X\) in the sense of
\cref{def:p03-morley-connected}: the permutations of any tuple of \(m\)
distinct elements include all \(\bar x^s\) for the increasing
arrangement \(\bar x\), some satisfying \(\psi\) and some not. The
Morleyization of \(T'\) is totally transcendental and \(X\) is infinite,
so \cref{thm:p03-morley-no-order} excludes this. Hence \(P\) is empty or
everything. Every tuple \(\bar u\) of \(m\) distinct elements is
\(\bar x^s\) for its increasing arrangement and some \(s\), and
\(D\models\psi(\bar u)\) exactly when \(P\) is everything. So any two such
tuples agree on \(\psi\).

\emph{The term hull and its automorphisms.} Let \(H=M(X,D)\), the set of
values in \(D\) of \(L^\#\)-terms at tuples from \(X\)
(\cref{def:p03-morley-skolem-hull}). It is closed under the function
symbols, hence a substructure of \(D\). The axioms of \(T^\#\) are
universal and hold in \(D\), so they hold in \(H\), and the reduct \(A\)
of \(H\) to the language of \(T\) is a model of \(T\). By
\cref{thm:p03-morley-ramsey-em}, \(\kappa(H)=\kappa(X)=\kappa\) since
\(X\) is infinite and \(L^\#\) countable; so \(\kappa(A)=\kappa(X)=\kappa\)
and \(X\subseteq|A|\).

Let \(\sigma\) be a one-one map of \(X\) onto itself, and set
\(\hat\sigma(t^D(\bar u))=t^D(\sigma(\bar u))\) for terms \(t\) and tuples
\(\bar u\) from \(X\). Any finite statement about elements of \(H\),
whether an equation \(t(\bar u)=t'(\bar u')\) or a relation
\(R(t_1(\bar u),\dots,t_n(\bar u))\) of \(L^\#\), can be written as
\(\chi(\bar z)\), where \(\bar z\) lists the distinct elements occurring
in the argument tuples and \(\chi\) is an \(L^\#\)-formula. The tuple
\(\sigma(\bar z)\) again consists of distinct elements, so by total
indiscernibility \(D\models\chi(\bar z)\) exactly when
\(D\models\chi(\sigma(\bar z))\). Applied to equations, this shows that
\(\hat\sigma\) is well defined. The map built in the same way from
\(\sigma^{-1}\) is its inverse, so \(\hat\sigma\) is a bijection of \(H\).
Applied to relations, it shows that \(\hat\sigma\) preserves and reflects
every relation symbol of \(L^\#\), in particular every relation symbol of
\(T\). Hence \(\hat\sigma\) is an automorphism of \(A\), and it extends
\(\sigma\) because each \(x\in X\) is the value of the term \(w_0\) at
\(x\).
\end{proof}
\end{annexitem}

\begin{annexitem}{annex:p03-w03-p03-033-jonsson-conditions}{thm:p03-w03-p03-033-jonsson-conditions}{Jónsson's conditions for the parameter systems}
\begin{proof}
\emph{Renaming.} Let \(A\) be a system of the similarity type of \(T\),
\(M\) a model of \(T\), and \(f:A\to M\) a monomorphism. Then there are a
model \(M'\) of \(T\) having \(A\) as a subsystem and an isomorphism
\(h:M\to M'\) with \(h(f(a))=a\) for every element \(a\) of \(A\). Indeed,
let \(Z=|M|\setminus f(|A|)\) and \(Z'=Z\times\{|A|\}\). No element
\((z,|A|)\) of \(Z'\) lies in \(|A|\): with Kuratowski pairs,
\(|A|\in\{z,|A|\}\in(z,|A|)\), and \((z,|A|)\in|A|\) would give a
membership cycle, which Foundation excludes. Put \(h=f^{-1}\) on
\(f(|A|)\) and \(h(z)=(z,|A|)\) for \(z\in Z\); this is a bijection of
\(|M|\) onto \(|A|\cup Z'\). Give \(|A|\cup Z'\) the relations
transported by \(h\), so that \(h\) is an isomorphism onto the resulting
system \(M'\), which is therefore a model of \(T\). For a relation symbol
\(R\) and a tuple \(\bar a\) from \(A\), \(R^{M'}(\bar a)\) holds exactly
when \(R^M(f(\bar a))\) holds, that is, exactly when \(R^A(\bar a)\)
holds, because \(f\) preserves and reflects \(R\). So \(A\) is the
subsystem of \(M'\) on \(|A|\).

\emph{Subsystems.} If \(B\in\mathcal N(T)\) is a subsystem of a model
\(M\), every subsystem of \(B\) is a subsystem of \(M\), hence belongs to
\(\mathcal N(T)\). In particular VI\(_\alpha\) holds for every ordinal
\(\alpha\) with \(C=B\): if \(A\in\mathcal N(T)\), \(B\prec A\) and
\(^*B<\aleph_\alpha\), then \(B\in\mathcal N(T)\) and
\(B\prec B\prec A\).

\emph{Condition II.} Let \(A\in\mathcal N(T)\) be a subsystem of a model
\(M\), and let \(k:A\to A'\) be an isomorphism. The map
\(k^{-1}:A'\to M\) is a monomorphism, and renaming gives a model having
\(A'\) as a subsystem. So \(A'\in\mathcal N(T)\).

\emph{Conditions I and I\('\).} \(T\) has an infinite model \(M\). Its
subsystems on a one-element and on a two-element subset of \(|M|\) belong
to \(\mathcal N(T)\) and are not isomorphic, since they have different
powers. By \cref{thm:upward-ls}, applied to \(T\) in its countable
language, \(T\) has a model of power \(\aleph_\xi\) for every ordinal
\(\xi\), and that model belongs to \(\mathcal N(T)\).

\emph{Conditions III and V.} These are the joint embedding clause and the
union clause of \cref{thm:p03-morley-normalization}; the union clause is
stated there for increasing families, which include the chains of
condition V.

\emph{Condition IV\('\).} Let \(f_0:A\to B_0\) and \(f_1:A\to B_1\) be
monomorphisms in \(\mathcal N(T)\), and let \(B_i\) be a subsystem of a
model \(M_i\). Renaming applied to \(f_i:A\to M_i\) gives models
\(M_i'\supseteq A\) and isomorphisms \(h_i:M_i\to M_i'\) with
\(h_if_i(x)=x\) for every element \(x\) of \(A\). Put
\(B_i'=h_i(B_i)\), the subsystem of \(M_i'\) on \(h_i(|B_i|)\); it
contains \(|A|\), so \(A\) is a common subsystem of \(B_0'\) and
\(B_1'\). The amalgamation clause of \cref{thm:p03-morley-normalization}
gives \(C\in\mathcal N(T)\) and monomorphisms \(e_i:B_i'\to C\) that agree
on \(A\). Then \(g_i=e_ih_i\) is a monomorphism of \(B_i\) into \(C\), and
\(g_0f_0(x)=e_0(x)=e_1(x)=g_1f_1(x)\) for every element \(x\) of \(A\).

\emph{Directed unions.} Let \(\mathcal D\) be a nonempty family of members
of \(\mathcal N(T)\) such that any two members are subsystems of a third
member. Then any two members induce the same relations on their common
elements, since both are subsystems of a third; so the union \(U\) of the
universes carries well-defined relations, and every member is a subsystem
of the system \(U\). Name each element of \(U\) by a new constant and let
\(\Delta\) be the set of atomic sentences and negated atomic sentences,
including equations and inequations of the constants, true in \(U\). A
finite subset of \(T\cup\Delta\) mentions finitely many elements of \(U\);
each lies in some member of \(\mathcal D\), and finitely many members have
a common upper bound \(D\in\mathcal D\). The finite subset is true in a
model of \(T\) having \(D\) as a subsystem, with each constant
interpreted by itself. By \cref{thm:p03-compactness-ultraproduct}, some
structure \(N\) satisfies \(T\cup\Delta\). Sending each element of \(U\)
to the interpretation of its constant is one-one, by the inequations in
\(\Delta\), and preserves and reflects every relation, by the atomic and
negated atomic sentences in \(\Delta\). Renaming applied to this
monomorphism into the reduct of \(N\), which is a model of \(T\), shows
that \(U\) is a subsystem of a model, so \(U\in\mathcal N(T)\).
\end{proof}
\end{annexitem}

\begin{annexitem}{annex:p03-w03-p03-033-jonsson-iv}{thm:p03-w03-p03-033-jonsson-iv}{Failure of condition IV for the integers}
\begin{proof}
The systems \(A\) and \(B\) are isomorphic to the subsystem of the
expansion of \((\mathbb Z,<)\) on \(\{0,1\}\), so they belong to
\(\mathcal N(T)\) by condition II
(\cref{thm:p03-w03-p03-033-jonsson-conditions}). The isomorphisms fix
\(0\), so the subsystems of \(A\) and of \(B\) on \(\{0\}\) are both
equal to the subsystem of the expansion of \((\mathbb Z,<)\) on
\(\{0\}\). Hence \(A\cap B\), whose universe is \(\{0\}\) because
\(a\neq b\) and neither is \(0\), is that one-element system; it belongs
to \(\mathcal N(T)\) and is a subsystem of \(A\) and of \(B\).

Let \(\sigma(v_0,v_1)\) be the formula
\(v_0<v_1\wedge\neg\exists v_2\,(v_0<v_2\wedge v_2<v_1)\), and let \(S\)
be the relation symbol of \(T\) for \(\sigma\). In \(\mathbb Z\),
\(\sigma(0,1)\) holds, so \(S^A(0,a)\) and \(S^B(0,b)\) hold. The
sentence
\(\forall x\,\forall y\,\forall z\,(\sigma(x,y)\wedge\sigma(x,z)\to y=z)\)
is true in \(\mathbb Z\), since an integer has at most one immediate
successor, so it belongs to \(\Sigma\), and its translation with \(S\) in
place of \(\sigma\) holds in every model of \(T\).

Suppose \(C\in\mathcal N(T)\) had \(A\) and \(B\) as subsystems, and let
\(C\) be a subsystem of a model \(M\) of \(T\). Then \(S^M(0,a)\) and
\(S^M(0,b)\) hold, and the sentence above gives \(a=b\), contrary to the
choice of \(a\) and \(b\). So no such \(C\) exists, and condition IV fails
for \(\mathcal N(T)\). Condition IV\('\) holds by
\cref{thm:p03-w03-p03-033-jonsson-conditions}; it is satisfied here by
identifying \(a\) with \(b\).
\end{proof}
\end{annexitem}

\begin{annexitem}{annex:p03-w03-p03-033-jonsson-algebraic}{thm:p03-w03-p03-033-jonsson-algebraic}{Algebraic points and algebraic elements}
\begin{proof}
For a formula \(\varphi\in F(A)\) and \(n<\omega\) write
\(N_n(\varphi)\) for the sentence
\(\exists v_0\cdots\exists v_n\bigl(\bigwedge_{i\leq n}\varphi(v_i)
\wedge\bigwedge_{i<j\leq n}v_i\neq v_j\bigr)\), which says that
\(\varphi\) has at least \(n+1\) solutions. Clause~4 says that
\(T(A)\vdash\neg N_n(\varphi)\) for some \(\varphi\in p\) and some \(n\).
Two facts are used.

\emph{(a) Models over \(A\).} If \(M\) is a model of \(T\) and
\(f:A\to M\) is a monomorphism, the expansion of \(M\) interpreting each
\(a\in A\) by \(f(a)\) is a model of \(T(A)\): it satisfies \(T\), and it
satisfies the diagram of \(A\) because \(f\) is one-one and preserves and
reflects every relation. Every \(C\supseteq A\) in \(\mathcal N(T)\) is a
subsystem of such a model, with \(f\) the inclusion.

\emph{(b) Equivalence is equality of points.} For \(C\supseteq A\) in
\(\mathcal N(T)\) and \(v\in C\), the map fixing \(A\) and sending \(u\)
to \(v\) is a monomorphism of \(A(u)\) onto \(A(v)\) exactly when \(u\)
and \(v\) realize the same point of \(S(A)\)
(\cref{thm:p03-morley-types}). Any monomorphism witnessing
\(A\)-equivalence is this map. So the elements of \(C\) that are
\(A\)-equivalent to \(u\) are exactly the elements of \(C\) that realize
\(p\).

\emph{A model with infinitely many realizations.} Suppose clause~4 fails.
Since \(T(A)\) is complete, \(T(A)\vdash N_n(\varphi)\) for every
\(\varphi\in p\) and every \(n\). Let \(c_i\), \(i<\omega\), be new
constants and let
\[
 \Gamma=T(A)\cup\{\varphi(c_i):\varphi\in p,\ i<\omega\}
 \cup\{c_i\neq c_j:i<j<\omega\}.
\]
A finite subset of \(\Gamma\) mentions finitely many formulas
\(\varphi_1,\dots,\varphi_m\) of \(p\) and constants among
\(c_0,\dots,c_k\). The conjunction \(\varphi\) of the \(\varphi_l\)
belongs to \(p\), and \(T(A)\vdash N_k(\varphi)\), so in any model of
\(T(A)\) the constants \(c_0,\dots,c_k\) can be interpreted by distinct
solutions of \(\varphi\). By \cref{thm:p03-compactness-ultraproduct},
\(\Gamma\) has a model \(N\). The map sending \(a\in A\) to the
interpretation of its constant is a monomorphism of \(A\) into the reduct
of \(N\), by the diagram of \(A\) in \(T(A)\); renaming
(the first step in the proof of \cref{thm:p03-w03-p03-033-jonsson-conditions}) gives a model
\(M\supseteq A\) of \(T\) isomorphic to that reduct over \(A\). The
images \(e_i\) of the interpretations of \(c_i\) are distinct, and each
satisfies every formula of \(p\); since \(p\) is a maximal consistent set
of formulas of \(F(A)\), each \(e_i\) realizes \(p\).

\emph{(4) implies (3).} Let \(\varphi\in p\) and
\(T(A)\vdash\neg N_n(\varphi)\). Let \(C\supseteq A\) in
\(\mathcal N(T)\), a subsystem of a model \(M\). By (a), \(M\) with the
elements of \(A\) named satisfies \(T(A)\), so \(\varphi\) has at most
\(n\) solutions in \(M\). By (b), the elements of \(C\) that are
\(A\)-equivalent to \(u\) realize \(p\), so they satisfy \(\varphi\);
there are at most \(n\) of them.

\emph{(3) implies (2)} is immediate.

\emph{(2) implies (4).} If clause~4 fails, the model \(M\supseteq A\)
constructed above has infinitely many elements \(e_i\) realizing \(p\).
By (b) they are \(A\)-equivalent to \(u\), and \(M\in\mathcal N(T)\), so
\(u\) is not algebraic over \(A\).

\emph{(4) implies (1).} Let \(\varphi\in p\) with
\(T(A)\vdash\neg N_n(\varphi)\). Recall \(S^0(A)=S(A)\), so
\(p\in S^0(A)\). Let \(f:A\to B'\) be a monomorphism in
\(\mathcal N(T)\) and let \(q\in(f^*)^{-1}(p)\). Since
\(f^*(q)=p\) contains \(\varphi\), the formula \(f(\varphi)\), obtained by
substituting \(f(a)\) for each parameter \(a\), belongs to \(q\). By
\cref{thm:p03-morley-types} there is an extension of \(B'\), which may be
taken to be a model \(M\) of \(T\), realizing every point of \(S(B')\).
Choose a realization \(e_r\in M\) of each \(r\in U_{f(\varphi)}\). Distinct
points have distinct realizations, since an element realizes only one
point, and every \(e_r\) satisfies \(f(\varphi)\) in \(M\). By (a) applied
to the monomorphism \(f:A\to M\), the formula \(f(\varphi)\) has at most
\(n\) solutions in \(M\). Hence \(U_{f(\varphi)}\) is a finite open subset
of the Hausdorff space \(S(B')\), and each of its points is isolated in
\(S(B')\): the point is the intersection of \(U_{f(\varphi)}\) with the
complement of the finitely many other points, which is open. In
particular \(q\) is isolated in \(S^0(B')=S(B')\). As \(f\) and \(q\) were
arbitrary, \(p\in\operatorname{Tr}^0(A)\) by
\cref{def:p03-morley-global-rank}.

\emph{(1) implies (4).} Suppose clause~4 fails and let \(M\supseteq A\)
and \(e_i\) be as constructed above, with \(i:A\to M\) the inclusion. Let
\(r_i\in S(M)\) be the point realized by \(e_i\). It contains
\(v_0=e_i\), so the \(r_i\) are distinct, and
\(i^*(r_i)=r_i\cap F(A)=p\). The fibre \((i^*)^{-1}(p)\) is closed in
\(S(M)\), as the inverse image of a point of the Hausdorff space \(S(A)\)
under a continuous map. The infinite set \(\{r_i:i<\omega\}\) of the
compact space \(S(M)\) has an accumulation point \(r\): otherwise each
point of \(S(M)\) has an open neighbourhood containing only finitely many
\(r_i\), and finitely many of these neighbourhoods cover \(S(M)\), which
is impossible. Since the fibre is closed, \(r\) lies in it, and \(r\) is
not isolated in \(S(M)=S^0(M)\). So the inclusion \(i\) has a fibre over
\(p\) with a non-isolated point, and \(p\notin\operatorname{Tr}^0(A)\).

\emph{The consequence.} \(B\) is an algebraic extension of \(A\) exactly
when every \(u\in B\) is algebraic over \(A\), which by the equivalence of
clauses 1 and 2 means that the point realized by every \(u\in B\) lies in
\(\operatorname{Tr}^0(A)\).
\end{proof}
\end{annexitem}

\begin{annexitem}{annex:p03-w03-p03-033-countable-prime}{thm:p03-w03-p03-033-countable-prime}{Prime models over a countable system}
\begin{proof}
\emph{Clause 1.} The theory \(T(A)\) is complete, its language \(L(A)\) is
countable because \(L\) and \(A\) are, and its models are models of \(T\),
which are infinite. By \cref{thm:p03-vaught-isolation-density}, from the
atomistic clause, \(T(A)\) has a prime model \(M_0\), and by
\cref{thm:p03-vaught-prime-atomic} it is countable; prime means that
\(M_0\) has an elementary embedding into every model of \(T(A)\)
(\cref{def:p03-vaught-model-notions}). The map sending \(a\in A\) to the
interpretation of its constant in \(M_0\) is a monomorphism of \(A\) into
the reduct of \(M_0\) to \(L\), by the diagram of \(A\) in \(T(A)\).
Renaming (the first step in the proof of \cref{thm:p03-w03-p03-033-jonsson-conditions}) gives a model
\(B\supseteq A\) of \(T\) and an isomorphism of the reduct of \(M_0\)
onto \(B\) carrying the interpretation of each constant \(a\) to \(a\).
With each \(a\in A\) naming itself, \(B\) is isomorphic to \(M_0\) as an
\(L(A)\)-structure, so it is a countable prime model of \(T(A)\).

Let \(C\) be a model of \(T\) and \(f:A\to C\) a monomorphism. The
expansion of \(C\) interpreting each \(a\in A\) by \(f(a)\) satisfies
\(T\) and the diagram of \(A\), because \(f\) is one-one and preserves and
reflects every relation, so it is a model of \(T(A)\). Primeness gives an
elementary embedding \(g\) of \(B\), with \(A\) named, into this
expansion. It sends each \(a\in A\) to the interpretation of its constant,
which is \(f(a)\), so \(g\) extends \(f\); and as an elementary embedding
it is one-one and preserves and reflects every relation of \(L\), so it is
a monomorphism of \(B\) into \(C\). Thus \(B\) is prime over \(A\).

\emph{Clause 2.} Assume that the isolated points are dense in \(S(A')\)
for every \(A'\in\mathcal N(T)\). We prove by induction on \(n\) that
every formula of \(F_n(T(A))\) consistent with \(T(A)\) has an atom below
it. For \(n=0\) the formulas are sentences; since \(T(A)\) is complete, a
consistent sentence \(\sigma\) satisfies \(T(A)\vdash\sigma\), so for
every sentence \(\gamma\) with \(\sigma\wedge\gamma\) consistent,
\(T(A)\vdash\gamma\) and hence \(T(A)\vdash\sigma\to\gamma\); thus
\(\sigma\) is an atom below itself.

Suppose the claim holds for \(n\), and let
\(\varphi(v_0,\dots,v_n)\) be consistent with \(T(A)\). The formula
\(\exists v_n\varphi\) of \(F_n(T(A))\) is consistent, so by the
induction hypothesis there is an atom \(\theta(v_0,\dots,v_{n-1})\) with
\(T(A)\vdash\theta\to\exists v_n\varphi\). Choose a model of \(T(A)\) and
a tuple \(\bar a=(a_0,\dots,a_{n-1})\) in it satisfying \(\theta\); after
renaming as in clause 1 we may take it to be a model \(M\supseteq A\) of
\(T\) in which each element of \(A\) names itself. Let \(A'\) be the
subsystem of \(M\) on \(|A|\cup\{a_0,\dots,a_{n-1}\}\). For a formula
\(\gamma(v_0,\dots,v_n)\) of \(L(A)\), let \(\gamma^{\bar a}\in F(A')\)
be the formula obtained by putting the constant for \(a_i\) in place of
\(v_i\) for \(i<n\) and then renaming \(v_n\) as \(v_0\). Every formula of
\(F(A')\) has this form: replace each constant for \(a_i\) by the
variable \(v_i\), using the least such \(i\) when several \(a_i\) are
equal, after renaming \(v_0\) as \(v_n\). By (a) in the proof of
\cref{thm:p03-w03-p03-033-jonsson-algebraic}, \(M\) with the elements of
\(A'\) named is a model of the complete theory \(T(A')\).

Since \(M\models\theta(\bar a)\), also \(M\models\exists v_n\varphi(\bar a,v_n)\),
so \(\varphi^{\bar a}\) is consistent with \(T(A')\) and
\(U_{\varphi^{\bar a}}\neq\varnothing\). By hypothesis \(U_{\varphi^{\bar a}}\)
contains an isolated point \(r\) of \(S(A')\). Choose a formula of
\(F(A')\) whose basic set is \(\{r\}\), and write it as \(\rho^{\bar a}\)
with \(\rho(v_0,\dots,v_n)\) a formula of \(L(A)\). Put
\(\chi=\theta\wedge\rho\in F_{n+1}(T(A))\).

\emph{Key step.} If \(\gamma(v_0,\dots,v_n)\) is a formula of \(L(A)\)
with \(\gamma^{\bar a}\in r\), then \(T(A)\vdash\chi\to\gamma\). Indeed,
\(U_{\rho^{\bar a}\wedge\neg\gamma^{\bar a}}=\{r\}\cap
U_{\neg\gamma^{\bar a}}=\varnothing\), so
\(\rho^{\bar a}\wedge\neg\gamma^{\bar a}\) is inconsistent with
\(T(A')\), since every consistent formula lies in some point of
\(S(A')\). Hence \(M\models\forall v_n\,(\rho(\bar a,v_n)\to\gamma(\bar a,v_n))\),
that is, \(\bar a\) satisfies the formula
\(\tau=\forall v_n(\rho\to\gamma)\) of \(F_n(T(A))\). Thus
\(\theta\wedge\tau\) is consistent, and since \(\theta\) is an atom,
\(T(A)\vdash\theta\to\tau\), which gives \(T(A)\vdash\chi\to\gamma\).

\emph{\(\chi\) lies below \(\varphi\).} Since \(r\in U_{\varphi^{\bar a}}\),
the key step with \(\gamma=\varphi\) gives \(T(A)\vdash\chi\to\varphi\).

\emph{\(\chi\) is consistent.} The point \(r\) contains \(\rho^{\bar a}\),
so \(\rho^{\bar a}\) is consistent with the complete theory \(T(A')\) and
\(\exists v_0\rho^{\bar a}\) belongs to it. Hence
\(M\models\exists v_n\rho(\bar a,v_n)\), and together with
\(M\models\theta(\bar a)\) this gives a tuple of \(M\) satisfying \(\chi\).

\emph{\(\chi\) is an atom.} Let \(\gamma\in F_{n+1}(T(A))\) with
\(\chi\wedge\gamma\) consistent. Either \(\gamma^{\bar a}\in r\) or
\((\neg\gamma)^{\bar a}\in r\). In the second case the key step gives
\(T(A)\vdash\chi\to\neg\gamma\), contradicting the consistency of
\(\chi\wedge\gamma\). So \(\gamma^{\bar a}\in r\), and the key step gives
\(T(A)\vdash\chi\to\gamma\).

This completes the induction, so every \(F_n(T(A))\) is atomistic.
\end{proof}
\end{annexitem}

\begin{annexitem}{annex:p03-w03-p03-033-coding-theory}{thm:p03-w03-p03-033-coding-theory}{The coding theory and the power \(\aleph_\omega\)}
\begin{proof}
\emph{The coding axioms.} For each \(k<\omega\) the sentence
\[
 \psi_k:\quad \forall y_0\cdots\forall y_k\,\exists x\,
 \bigwedge_{i\leq k}f_i(x)=y_i
\]
is true in \((Y,f_n)_{n<\omega}\): given \(y_0,\dots,y_k\in Y\), define
\(x\in Y\) by \(x(\langle i,m\rangle)=y_i(m)\) for \(i\leq k\) and
\(m<\omega\), and \(x(j)=0\) for every \(j\) not of the form
\(\langle i,m\rangle\) with \(i\leq k\); this is well defined because
\(\langle\cdot,\cdot\rangle\) is a bijection, and \(f_i(x)=y_i\) for
\(i\leq k\). So \(\psi_k\in\Sigma_{\mathrm{cd}}\), and in every model of
\(T_{\mathrm{cd}}\) the translation
\(\forall y_0\cdots\forall y_k\,\exists x\bigwedge_{i\leq k}G_i(x,y_i)\)
holds, as does the statement that each \(G_n\) is the graph of a total
function, by \cref{def:p03-morley-standing} and
\cref{thm:p03-morley-normalization}. The theory \(\Sigma_{\mathrm{cd}}\)
is complete, being the theory of one structure, its language is
countable, and \(Y\) is infinite.

\emph{Clause 1.} Let \(a_0,a_1\in Y\) be the constant functions with
values \(0\) and \(1\), and let \(A\) be the subsystem of the expansion of
\(Y\) on \(\{a_0,a_1\}\). For \(s\in\{0,1\}^\omega\) define \(x_s\in Y\)
by \(x_s(\langle n,m\rangle)=a_{s(n)}(m)\); then \(f_n(x_s)=a_{s(n)}\) for
every \(n\). Let \(p_s\in S(A)\) be the point realized by \(x_s\); it
contains \(G_n(v_0,a_{s(n)})\) for every \(n\). If \(s\neq t\), choose
\(n\) with \(s(n)\neq t(n)\). Because \(G_n\) is the graph of a function
and \(a_0\neq a_1\), the formula
\(G_n(v_0,a_0)\wedge G_n(v_0,a_1)\) is inconsistent with \(T(A)\), so no
point contains both \(G_n(v_0,a_{s(n)})\) and \(G_n(v_0,a_{t(n)})\), and
\(p_s\neq p_t\). Hence \(S(A)\) has at least \(2^{\aleph_0}\) points. It
has at most \(2^{\aleph_0}\), since \(F(A)\) is countable and each point is
a subset of it. If \(T_{\mathrm{cd}}\) were totally transcendental,
\cref{thm:p03-morley-countability} would give
\(\kappa(S(A))\leq\kappa(A)+\aleph_0=\aleph_0\), which is false. So
\(T_{\mathrm{cd}}\) is not totally transcendental.

\emph{Clause 2.} Let \(B\) be a model of \(T_{\mathrm{cd}}\) of uncountable
power \(\lambda\) that is saturated over countable subsystems. For
\(s:\omega\to|B|\), let \(A_s\) be the subsystem of \(B\) on the range of
\(s\), which is countable, and let
\(Q_s=\{G_n(v_0,s(n)):n<\omega\}\subseteq F(A_s)\). For finitely many
formulas of \(Q_s\), with indices at most \(k\), the translated coding
axiom gives \(x\in|B|\) with \(G_i(x,s(i))\) for every \(i\leq k\). Since
\(B\) with the elements of \(A_s\) named is a model of \(T(A_s)\) (fact (a)
in the proof of \cref{thm:p03-w03-p03-033-jonsson-algebraic}), the
conjunction of these formulas is consistent with \(T(A_s)\), and its
basic set in \(S(A_s)\) is nonempty. The basic sets of the members of
\(Q_s\) are closed sets with the finite intersection property in the
compact space \(S(A_s)\) (\cref{thm:p03-morley-types}), so some point
\(p_s\in S(A_s)\) contains \(Q_s\). Since \(\kappa(A_s)\leq\aleph_0\),
saturation over countable subsystems gives \(b_s\in|B|\) realizing
\(p_s\), so \(G_n^B(b_s,s(n))\) holds for every \(n\). As \(G_n^B\) is the
graph of a function, \(s(n)\) is the unique \(y\) with \(G_n^B(b_s,y)\),
so \(s\) is determined by \(b_s\). Hence \(s\mapsto b_s\) is a one-one map
of \(|B|^\omega\) into \(|B|\), and
\(\lambda^{\aleph_0}\leq\lambda\leq\lambda^{\aleph_0}\).

\emph{Clause 3.} Let \(\kappa=\aleph_\omega\), the union of the increasing
sequence of ordinals \(\aleph_n\), \(n<\omega\). Suppose that
\(\xi\mapsto h_\xi\) maps \(\kappa\) onto the set \(\kappa^\omega\) of
functions from \(\omega\) to \(\kappa\). For \(n<\omega\) the set
\(\{h_\xi(n):\xi<\aleph_n\}\) has power at most \(\aleph_n<\kappa\), so
it omits some ordinal below \(\kappa\); let \(h(n)\) be the least such
ordinal. Then \(h\in\kappa^\omega\). For every \(\xi<\kappa\) there is
\(n\) with \(\xi<\aleph_n\), and then \(h(n)\neq h_\xi(n)\), so
\(h\neq h_\xi\). This contradicts the choice of \(\xi\mapsto h_\xi\). Hence
\(\kappa^{\aleph_0}\), the power of \(\kappa^\omega\), is greater than
\(\kappa\). By \cref{thm:upward-ls}, \(\Sigma_{\mathrm{cd}}\) has a
model of power \(\aleph_\omega\), and its unique expansion
(\cref{thm:p03-morley-normalization}) is a model of \(T_{\mathrm{cd}}\) of
the same power. By clause~2, no model of \(T_{\mathrm{cd}}\) of power
\(\aleph_\omega\) is saturated over countable subsystems.

\emph{The consequence.} If \(\kappa>\aleph_0\) and
\(\kappa^{\aleph_0}=\kappa\), a model of power \(\kappa\) saturated over
countable subsystems exists by \cref{thm:p03-w03-p03-034-mv-existence},
which assumes no total transcendence. Conversely, such a model of power
\(\kappa\) forces \(\kappa^{\aleph_0}=\kappa\) by clause~2.
\end{proof}
\end{annexitem}
